
\section{Related Work}

\subsection{Supervised Optical Flow and the Shift Toward Global Correspondence}

Supervised optical flow has traditionally been driven by architectures trained on large synthetic datasets with dense ground-truth motion. Early end-to-end CNN models such as FlowNet and FlowNet2 \cite{dosovitskiy2015flownet,ilg2017flownet2} first demonstrated that dense correspondence could be learned directly from image pairs. Later methods, including SpyNet, PWC-Net \cite{pwcnet}, and LiteFlowNet \cite{hui2018liteflownet}, improved efficiency and estimation quality through pyramid processing, warping, and local cost volumes. Although effective, these models remain largely dependent on local correlation windows and hand-designed coarse-to-fine heuristics, which can be limiting in large-displacement or weak-texture scenarios.

A major advance came with RAFT \cite{raft}, which replaced coarse-to-fine estimation with all-pairs correlation and recurrent iterative refinement. This paradigm substantially improved accuracy and inspired many follow-up methods such as GMA \cite{gma}, CRAFT \cite{craft}, and SKFlow \cite{skflow}. More recently, Transformer-based approaches such as GMFlow \cite{gmflow}, FlowFormer, and FlowFormer++ \cite{flowformer,flowformerpp} have emphasized global correspondence modeling through attention or latent cost tokens, enabling stronger long-range reasoning than earlier CNN pipelines. In parallel, supervised multi-frame methods such as VideoFlow \cite{videoflow2023} show that temporal aggregation further improves robustness in occluded, ambiguous, and large-motion regions. These developments suggest that both global matching and temporal reasoning are important for robust flow estimation, especially in domains where local appearance cues are unreliable.

\subsection{Unsupervised and Self-Supervised Optical Flow}

Unsupervised optical flow aims to remove the need for dense annotations by learning from raw videos with reconstruction-based objectives. Early approaches typically rely on photometric consistency between a target frame and a warped source frame, together with edge-aware smoothness and forward-backward consistency constraints. Representative methods include UnFlow \cite{unflow}, DDFlow \cite{ddflow}, UPFlow \cite{upflow}, and ASFlow \cite{asflow}, which progressively improved occlusion handling, data distillation, pyramid refinement, and training stability. However, most of these methods still operate on two-frame inputs, making it difficult to exploit longer temporal context or maintain consistency across extended sequences.

Recent work has therefore moved toward richer self-supervised strategies. SMURF \cite{smurf} demonstrates that self-teaching, sequence-aware supervision, and multi-frame training signals can substantially improve unsupervised flow learning while retaining two-frame inference. MDFlow \cite{mdflow} further improves unsupervised estimation by selecting more reliable pseudo-labels and reducing the effect of ambiguous motion regions. Related approaches also explore teacher-student learning and domain adaptation to improve generalization beyond the source training distribution. For example, FlowDA \cite{flowda} introduces a mean-teacher based unsupervised domain adaptation framework to reduce the gap between synthetic and real data. Overall, this line of work shows that modern unsupervised optical flow is increasingly shaped by temporal reasoning, self-distillation, and domain-aware learning rather than by photometric losses alone.

\subsection{Structure-Guided Motion Estimation for Animation and Stylized Video}

A growing body of work recognizes that higher-level structural priors can substantially improve motion estimation when low-level appearance cues are ambiguous. Segmentation masks, object boundaries, and region-level regularization can encourage coherent motion within objects while preserving discontinuities across contours. This is particularly relevant for animation, where object interiors are often nearly textureless and perceptually important motion cues concentrate around boundaries. Recent unsupervised methods such as UnSAMFlow \cite{unsamflow} show that segmentation priors from foundation models can be integrated directly into flow learning to improve motion boundaries, object coherence, and cross-domain generalization.

These ideas are especially important for 2D animation, a domain that remains underexplored compared with natural-image optical flow. The AnimeRun benchmark \cite{animerun} highlights the substantial domain gap between natural scenes and anime frames, where piecewise smooth regions, repeated colors, exaggerated deformation, and hard contours often violate assumptions underlying natural-video optical flow. At the application level, animation-specific systems such as AnimeInterp \cite{animeinterp} and line-art interpolation methods \cite{linedrawingFlow} further show the practical importance of accurate motion estimation for stylized content. Nevertheless, there is still limited work that jointly combines unsupervised learning, multi-frame temporal modeling, and structure-guided reasoning for animation optical flow. Our method is motivated by this gap and brings these three directions into a unified framework for 2D animation.

